\section{Research Applications}

The pilot results provide an initial indication of how structured lineage data may support the applications discussed below.

A mature TPR ecosystem would enable at least three substantive research programs.

First, it would provide a stronger empirical foundation for computational memetics and cultural evolution. Instead of relying primarily on analogy, toy models, or digital meme corpora, researchers could test transmission and retention models against structured histories of material artifacts.

Second, it would support richer convergence analysis and innovation forecasting. Patent-only systems can detect some signals, but they miss informal redesign and consumer-driven pressure [15, 16, 17, 18, 19, 20, 23, 37, 39]. A repository that captures both lineage structure and broader selection signals could produce more informative convergence maps.

Third, it would support a more rigorous study of multiple discovery. Independent invention is well documented, but we still understand relatively little about the pressure configurations under which it becomes likely. TPRs could support event comparison at a level of detail not currently available, especially when linked to cross-domain evolution pathways and convergence signals [20, 23, 37, 39].
